\documentclass[
    12pt,
    openright,
    % twoside,
    a4paper,
    brazil
]{abntex2}

% Pacotres básicos
\usepackage{lmodern}
\usepackage[T1]{fontenc}
\usepackage[utf8]{inputenc}
\usepackage{indentfirst}
\usepackage{color}
\usepackage{graphicx}
\usepackage{hyperref}
\usepackage{microtype}
\usepackage[alf]{abntex2cite}

\titulo{Material de Apoio -- Seminário de Prolog}

\autor{
    Amanda Andrade Abreu 
    \\
    Guilherme Amaral Giffoni
}
\local{Belo Horizonte}
\data{2026}
\orientador[Professor:]{Dr. Marco Rodrigo Costa}
\instituicao{
    PONTIFÍCIA UNIVERSIDADE CATÓLICA DE MINAS GERAIS
    \par
    Instituto de Ciências Exatas e Informática
}
\tipotrabalho{Material de Apoio}

\definecolor{blue}{RGB}{41,5,195}

\makeatletter
\hypersetup{
    pdftitle={\@title},
    pdfauthor={\@author},
    pdfsubject={\imprimirpreambulo},
    pdfcreator={LaTeX with abnTeX2},
    colorlinks=true,
    citecolor=blue,
    filecolor=magenta,
    urlcolor=blue
    urlcolor=blue
}

\begin{document}

\imprimirfolhaderosto*

% SUMÁRIO
\pdfbookmark[0]{\contentsname}{toc}
\tableofcontents*
\cleardoublepage

% falta roteiro e agenda...

% TEXTO PRINCIPAL
\textual
\chapter{Introdução}

\chapter{Histórico}
% cronologia e genealogia

\chapter{Paradigmas da Linguagem}

\chapter{Características Marcantes}

\chapter{Linguagens Relacionadas}
% influenciadores, influenciadas, similares, "opostas", etc

\chapter{Exemplo de Programas em Prolog}

\chapter{Considerações Finais}

\chapter{Bibliografia}

\end{document}